\documentclass[12pt]{article}
% Préambule commun à tous les documents extraits de « La Revue CAPES »
\usepackage[T1]{fontenc}
\usepackage[utf8]{inputenc}
\usepackage{lmodern}
\usepackage{amsmath,amssymb,amsthm,amsfonts,mathrsfs}
\usepackage{setspace}
\usepackage{listings}
\usepackage{pgfplots}
\pgfplotsset{compat=1.18}
\usepackage{longtable}
\usepackage{adjustbox}
\usepackage{lettrine}
\usepackage{tkz-tab}
\usepackage{changepage}
\usepackage{pdflscape}
\usepackage{graphicx}
\usepackage{tikz}
\usepackage{xcolor}
\usepackage[a4paper,top=2cm,bottom=2.5cm,left=2.5cm,right=2cm]{geometry}
\usepackage{fancyhdr}
\usepackage{draftwatermark}
\SetWatermarkText{La RevueCapes}
\SetWatermarkLightness{0.9}
\SetWatermarkScale{2}
\usepackage{hyperref}
\hypersetup{colorlinks=true,linkcolor=blue!50!black,urlcolor=blue!50!black}

\definecolor{revue}{RGB}{20,60,120}

\pagestyle{fancy}
\fancyhf{}
\renewcommand{\headrulewidth}{0.4pt}
\setlength{\headheight}{15pt}

% \entete{Matière}{Titre}{Type de document}
\newcommand{\entete}[3]{%
  \fancyhead[L]{\small\textsc{La Revue CAPES} -- #1}%
  \fancyhead[R]{\small #2}%
  \fancyfoot[C]{\thepage}%
  \thispagestyle{plain}%
  \begin{center}
    {\color{revue}\rule{\textwidth}{1.2pt}}\\[4pt]
    {\small\textsc{Concours directs CAP-CEG / CAPES -- Burkina Faso}}\\[6pt]
    {\LARGE\bfseries\color{revue} #2}\\[6pt]
    {\large\itshape #1 -- #3}\\[2pt]
    {\color{revue}\rule{\textwidth}{1.2pt}}
  \end{center}
  \vspace{0.5cm}%
}

% Encadré pour les remarques ajoutées dans les corrigés
\newcommand{\remarque}[1]{\par\medskip\noindent\fcolorbox{revue}{revue!6}{\parbox{\dimexpr\linewidth-2\fboxsep-2\fboxrule}{\textbf{\color{revue}Remarque.} #1}}\par\medskip}


\begin{document}
\entete{Mathématiques}{Ancien sujet CAPES -- Session 2022}{Énoncé}
\begin{spacing}{1.5}

\textbf{Exercice 1 : (4pts)}

Soit la suite numérique $(U_n)$ définie sur $\mathbb{N}$ par : $U_0 = 2$ et $U_{n+1} = \frac{2}{3}U_n + \frac{1}{3}n + 1, \forall {n\in \mathbb{N}}$.
1.a. Démontrer que pour tout $n, U_n\leq n+3$.

b. Démontrer que $\forall {n\in \mathbb{N}}, U_{n+1} - U_n = \frac{1}{3}(n+3-U_n)$.

c. En déduire le sens de variation de la suite $(U_n)$.

2. On désigne par $(V_n)$ la suite définie sur $\mathbb{N}$ par $V_n =U_n-n$.

a. Démontrer que la suite $(V_n)$ est une suite géométrique de raison $\frac{2}{3}$.

b. En déduire $\forall {n\in \mathbb{N}}$ que $U_n = 2(\frac{2}{3})^n + n$.

c. La suite $(U_n)$ est-elle convergente?

\textbf{Exercice 2 : (6pts)}

Soit l'intégrale $I=\int_{0}^{1}\frac{\ln(1 + x)}{1 + x^2}dx$.

1. Montrer que $\forall {x\in ]-1,+\infty[}$, on a : $\ln(1+x) = \int_{0}^{1}\frac{x}{1 + xy}dy$.

En déduire que $I = \int\int_D\frac{x}{(1+x^2)(1+xy)}dxdy$ ou  $D$ est le pavé $[0,1]^2$.

2. Considérons $f(x)= \frac{x}{(1+x^2)(1+xy)}$ pour $y$ fixé.

Décomposer $f(x)$ en élément simple puis calculer $\int_{0}^{1}f(x)dx$.

Montrer que $2I = \int_{0}^{1}\frac{1}{1+y^2}(\frac{1}{2}ln2 + y\frac{\pi}{4}ln2)dy$.

En déduire que $I=\frac{\pi}{8}ln(2)$

\medskip
\textbf{Exercice 3 : (5pts)}

On considère le système d'équations différentielle suivant :  $(S_1)\begin{cases}
x^{\prime}(t)=-x(t) + 3 y(t)\\
y^{\prime}(t)=x(t) + y(t)
\end{cases}$

1. Montrer que $(S_1)\Longleftrightarrow X^{\prime} = AX$ ou $A$ et $X$ sont à déterminer.

2. Montrer que $A$ est diagonalisable et l'écrire sous la forme $A=PDP^{-1}$. On précisera $D; P$ et $P^{-1}$.

3. Montrer que $(S_1)\Longleftrightarrow(S_2):Y^{\prime} = DY$. Préciser $Y$.

4. Résoudre $(S_2)$ puis déduire la solution générale de $(S_1)$.

\textbf{Exercice 4 : (5pts}

On considère deux fonctions $f$ et $g$ définies sur $[\frac{3\pi}{2},2\pi]$ par :

$f(x) = \cos x$ et $g(x)=\frac{\cos x}{1 - sin x}$.

Soit $C_f$ et $C_g$ les courbes représentatives de $f$ et $g$ respectivement dans un repère orthonormé $(o;\overrightarrow{i};\overrightarrow{j})$.

1. On suppose que pour tout $x\in [\frac{3\pi}{2},2\pi], g(x) - f(x) = \frac{h(x)}{1-\sin x}$.

Donner une expression simplifiée de $h(x)$.

2. Déterminer l'ensemble $S$ des solutions de l'équation $g(x)-f(x)=0$ dans $[\frac{3\pi}{2},2\pi]$.

3. Étudier le signe de $g(x)-f(x)$ sur $[\frac{3\pi}{2},2\pi]$. Préciser les positions relatives des courbes $C_f$ et $C_g$.

4.a. Déterminer $g^{\prime}(x)$ ou $g^{\prime}$ est la dérivée de $g$.

b. Dresser le tableau de variations de $g$.












\newpage

\end{spacing}
\end{document}
